%%%%%%%%%%%%%%%%%%%%%%%%%%%%%%%%%%%%%%%%%%%%%%%%%%%%%%%%%%%%%%
%%%%%%%%%%%%% MA-0551 Principios de Análisis en varias variables %%%%%%%%%%%%%%%%
%%%%%%%%%%%%%%%%%%%%%%%%%%%%%%%%%%%%%%%%%%%%%%%%%%%%%%%%%%%%%%
\newpage
\subsection*{MA-0551 Principios de Análisis II}
\addcontentsline{toc}{subsection}{MA-0551 Principios de Análisis II}

\hypertarget{MA0551}{}
\begin{refsection}
\begin{table}[H]
\centering{
\begin{tabular}{|ll|}
\hline
%\multicolumn{2}{|c|}{\textbf{Nivel de virtualidad del curso:} Virtual}\\ \hline
\textbf{Año:} III  & \textbf{Requisitos:} MA-0451 Principios de Análisis I y \\ & MA-0461 Álgebra Lineal II \\ 
\textbf{Ciclo:} V  & \textbf{Correquisitos:} Ninguno \\ 
\textbf{Tipo de curso:} Teórico & \textbf{Horas presenciales por semana:}   \\ 
\textbf{Créditos:} 4 & \textbf{Horas de trabajo independiente por semana:}  \\ \hline
\end{tabular}
}
\end{table}



\subsubsection*{Descripción del curso}

% Integración de Riemann en R
% Funciones continuas en R^n
% Diferenciación en R^n
% Teorema de la función inversa
% Teoremas de Stokes


Este es el segundo de una secuencia de dos cursos que introducen los principios básicos del análisis matemático en una y varias variables reales. A lo largo de esta secuencia se cubren los temas usuales del análisis matemático
de la recta real y el cálculo diferencial e integral en una y varias variables. 

El curso se encuentra enfocado al desarrollo del formalismo y la abstracción de elementos de análisis matemático en en varias variables reales, mediante la utilización de técnicas de razonamiento y demostración propias del área.

Se continúa con la formalización de las nociones y propiedades básicas del cálculo diferencial e integral, extendiéndolas a varias variables reales. Además, como apoyo en esta formalización, se toman las intuiciones y destrezas operacionales desarrolladas en los cursos MA-02XX Introducción al cálculo en una variable y MA-03XX Introducción al cálculo en varias variables, y los fundamentos de lógica y estrategias de demostración adquiridas en los cursos MA-01XX Matemática Exploratoria y MA-02XX Introducción a las Demostraciones.





\subsubsection*{Objetivos}

\textbf{General:} Aplicar los conceptos básicos de convergencia, diferenciación e integración para la demostración de propiedades del análisis en una y varias variables reales. \\

\textbf{Específicos:}

\begin{enumerate}
\item Asociar los conceptos y propiedades de convergencia, continuidad, diferenciación e integración en una variable con los correspondientes en varias variables.
\item Demostrar resultados sobre límites, continuidad, convergencia, diferenciación e integración en el análisis de una y varias variables reales.
\item Aplicar los teoremas de diferenciación e integración en varias variables en la demostración de resultados y la resolución de problemas teóricos y prácticos.
\item Indagar en la literatura sobre temas que permitan complementar su formación.
\item Comunicar ideas matemáticas de manera clara y efectiva en forma oral y escrita.
\end{enumerate}

\subsubsection*{Contenidos}

\begin{enumerate}
\item \textbf{Integración de Riemann en $\mathbb{R}$ (3 semanas):} 
\begin{enumerate}
    \item Particiones, refinamientos y propiedades básicas.
    \item Sumas de Darboux y sumas de Riemann.
    \item Integrabilidad según Riemann y según Darboux, y su equivalencia. Criterios de integrabilidad.
    \item Integrabilidad y continuidad.
    \item Propiedades de la integral de Riemann: linealidad, aditividad de intervalo, Teorema del Valor Medio, monotononía. 
    \item El Teorema Fundamental del Cálculo y aplicaciones. 
    \item Integración y convergencia uniforme.
    \item Integrales impropias, diferenciación bajo el signo de integral.
\end{enumerate}

\item \textbf{Topología métrica de} $\mathbb{R}^n$ \textbf{(3 semanas):}
\begin{enumerate}
\item Definiciones b\'asicas: Distancia, bolas, conjuntos abiertos y cerrados.
\item Sucesiones, convergencia, Teorema de Bolzano-Weierstrass.
\item Conjuntos conexos y propiedades de los mismos.
\item Conjuntos compactos, compacidad secuencial, Teorema de Heine-Borel, teoremas relacionados.
\end{enumerate}

\item \textbf{Funciones de varias variables (1 semana):}
\begin{enumerate}
%\item Dominio y rango de funciones de $\mathbb{R}^n$ en $\mathbb{R}$.  Funciones de $\mathbb{R}^n$ en $\mathbb{R}^m$.
\item Limites de funciones en 2 o 3 variables: C\'alculo del l\'imite, y test de caminos para probar la no existencia. Propiedades de los límites.
\item Continuidad: Definición formal, definición topológica, caracterización por sucesiones.  Verificación de la continuidad de funciones dadas.  Continuidad en conjuntos conexos y compactos.
\item Continuidad uniforme: definición, ejemplos y caracterizaciones.
\end{enumerate}


\item \textbf{Diferenciaci\'on en varias variables (4 semanas):}
\begin{enumerate}
\item Definición de función diferenciable de $\mathbb{R}^n$ en $\mathbb{R}^m$.  Derivada como transformación lineal.  Matriz Jacobiana. Propiedades de la derivada: linealidad, regla del producto (escalar y vectorial).  Relación entre diferenciabilidad y continuidad.
\item Regla de la cadena: Como composición de transformaciones lineales o producto de matrices Jacobianas.
\item Relación de las derivadas parciales con el concepto de diferencial.
\item Los teoremas de la función inversa y de la función implícita. Aplicaciones teóricas y prácticas.
\item Fórmula de Taylor en varias variables: Demostración del Teorema de Taylor, y cálculo de aproximaciones polinomiales de funciones de 2 o 3 variables.  Aplicaciones a la aproximación y cálculo de límites.
\item Optimización en varias variables: Puntos críticos y su clasificación en regiones abiertas por medio del criterio del Hessiano y segundas derivadas. Máximos y mínimos en regiones compactas.
\item Multiplicadores de Lagrange: Demostración del teorema, puntos críticos ligados con funciones de 2 o tres variables, con una o dos restricciones; clasificación de puntos críticos por medio de segundas derivadas y el Hessiano Orlado.
\end{enumerate}


% \item \textbf{Integraci\'on m\'ultiple:}
% \begin{enumerate}
% \item Construcci\'on de la integral sobre un rect\'angulo en $\mathbb{R}^n$ y propiedades elementales. 
% \item Contenido y conjuntos Jordan-medibles; definici\'on, propiedades, caracterizaciones.
% \item Integrales m\'ultiples sobre regiones m\'as generales.
% \item Teorema de Fubini y cambio de orden de integraci\'on.  C\'alculo de integrales dobles y triples sobre regiones generales.
% \item Teorema de cambio de variables: Demostraci\'on y c\'alculo de integrales por medio de sustituci\'on en varias variables.  Coordenadas polares, cil\'indricas y esf\'ericas.
% \end{enumerate}



\item \textbf{C\'alculo y an\'alisis vectorial (4 semanas):} 
\begin{enumerate}
\item Curvas en $\mathbb{R}^n$: Parametrizaci\'on por funciones vectoriales, suavidad, rectificabilidad, longitud de arco, derivaci\'on de funciones vectoriales, curvatura, vectores tangente, normal y binormal.
\item Integrales de l\'inea: Definici\'on para funciones escalares y campos vectoriales, propiedades y c\'alculo para curvas suaves por medio de una parametrizaci\'on.
\item Superficies: Definici\'on y ejemplos de superficies elementales (cu\'adricas), parametrizaci\'on, \'area de superficie.
\item Integrales de superficie: Definici\'on para superficies suaves y c\'alculo de integrales de superficies suaves por medio de parametrizaciones. 
\item Independencia de caminos: Campos vectoriales conservativos, funciones potenciales, Teorema Fundamental de Integrales de L\'inea, c\'alculo de integrales por medio de funciones potenciales. 
\item Teoremas de Green, Stokes y Gauss: Demostraci\'on de casos particulares, y c\'alculo de integrales utilizando dichos teoremas. 
\end{enumerate}

\end{enumerate}

\subsubsection*{Metodología}

\noindent \textbf{Lineamientos metodológicos:} La persona docente a cargo de este curso debe respetar las siguientes pautas:
\begin{enumerate}
    \item La presentación de los contenidos debe priorizar la formalización de los conceptos estudiados en los cursos MA-02XX Introducción al Cálculo en una Variable y MA-03XX Introducción al Cálculo en Varias Variables. En particular en este curso se debe asumir que las personas estudiantes  ya manejan las intuiciones y las habilidades operacionales sobre los conceptos básicos del cálculo diferencial e integral.
    \item En este curso se aplican las bases de lógica y las estrategias de demostración vistas en el curso MA-02XX Introducción a las Demostraciones.
    \item Las demostraciones deben presentarse tomando en cuenta que están dirigidas a una persona estudiante principiante. En particular, se debe priorizar dar un mayor nivel de detalle por encima de la brevedad y la estética, cuando esto ayude a mejorar la comprensión de los temas.
    \item Cuando la naturaleza del contenido lo permita, se deben utilizar representaciones visuales que apoyen la comprensión de los temas.
    \item Al presentar los teoremas de la función implícita y de la función inversa no deben hacerse las demostraciones pues estas no aportan al desarrollo de las habilidades planteadas en los objetivos del curso y además consumen una gran cantidad de tiempo. Sin embargo, sí es imprescindible que se muestre el uso de estos teoremas en aplicaciones teóricas y prácticas.
\end{enumerate}

\textbf{Sugerencias metodológicas:} Adicionalmente, se tienen las siguientes recomendaciones para la persona docente a cargo de este curso:
\begin{enumerate}
    \item Utilizar el recurso tecnológico para reforzar distintos conceptos estudiados en el curso por medio de la visualización, cálculo y experimentación. Se pone a disposición de la persona docente un repositorio de herramientas listas para ser implementadas en el curso, tanto en las clases como para el trabajo independiente de las personas estudiantes.
    \item Aprovechar momentos adecuados del curso para motivar el estudio por medio de reseñas acerca del desarrollo histórico del análisis matemático y su importancia en aplicaciones dentro y fuera de la matemática.
    \item En la evaluación, se sugiere utilizar estrategias que promuevan el trabajo constante de parte de las personas estudiantes a lo largo del ciclo lectivo. Por ejemplo, esto se puede hacer por medio de tareas o quizzes.
    \item Para fomentar el desarrollo de las habilidades investigativas en el estudiantado, se sugiere la implementación de pequeños proyectos de investigación. Por ejemplo, pueden ser proyectos de investigación bibliográfica o también proyectos cortos donde se desarrolle un tema por medio de ejercicios dirigidos, como se propone, por ejemplo, en el libro de Bartle \cite{Bar76}.
    \item Dado que hoy en día la mayoría del trabajo en matemática, tanto a nivel académico como profesional, es de naturaleza colaborativa, se sugiere a la persona docente implementar estrategias que promuevan el trabajo en equipo.
    \item Se sugiere hacer un recordatorio del Teorema del Cambio de Variable para integrales múltiples previo a empezar el estudio de integrales de superficie.
\end{enumerate}




\nocite{Abb15}
\nocite{Apo84}
\nocite{Apo85} 
\nocite{Bar76}
\nocite{BS11}
\nocite{Apo76}
%\nocite{SRW17}
%\nocite{Ste18}
%\nocite{Spi08}
\nocite{Ros13}
\nocite{TBB08}
\nocite{Rud76}
\nocite{Piz06}
\nocite{Buc03}
\nocite{Fit09}
\nocite{Spi65}
\nocite{MT03}

\printbibliography[heading=subbibliography]
\end{refsection}

